% ── Résumé (Français) ─────────────────────────────────────────
\chapter*{Résumé}
\addcontentsline{toc}{chapter}{Résumé}

Ce projet de fin d'études, réalisé au sein de l'entreprise \textbf{CIVCO BTP}, porte sur la
conception et le développement d'une application web dédiée à la gestion et à la planification
des projets de construction de l'entreprise. Face à la multiplication des chantiers et à la
diversité des intervenants, CIVCO BTP souhaitait disposer d'un outil numérique centralisé,
capable de structurer l'activité quotidienne et de faciliter le partage d'information entre
les services.

\vspace{0.35cm}

L'application couvre l'ensemble du cycle de vie d'un projet : création et suivi des chantiers
par phases, gestion des clients et de leurs contacts, affectation des équipes, planification
des tâches, établissement des devis et des factures, production des bons de livraison, suivi
des contrats et consultation d'un tableau de bord synthétique. Elle vise à remplacer les outils
dispersés (tableurs, documents papier, échanges informels) par une plateforme unique, accessible
aux différents profils métier de l'entreprise selon leurs droits d'accès.

\vspace{0.35cm}

La solution repose sur une architecture \textit{full-stack} : un \textit{backend} API
développé avec \textbf{Laravel} et \textbf{PHP}, exposant des services \textbf{REST} sécurisés
et s'appuyant sur une base de données relationnelle, couplé à une interface utilisateur
\textbf{React} moderne organisée en \textbf{SPA} (\textit{Single Page Application}). Un système
de rôles et de permissions (\textbf{RBAC}) garantit que chaque utilisateur accède uniquement
aux fonctionnalités correspondant à son profil.

\vspace{0.35cm}

Ce rapport présente le contexte du projet, les choix technologiques retenus, la conception
fonctionnelle et technique du système, ainsi que la réalisation et la validation des principaux
modules développés. Il s'articule autour de quatre chapitres : cadrage et méthodologie,
état de l'art, conception \textbf{UML} du système, puis implémentation détaillée des
fonctionnalités livrées.

\vspace{0.35cm}

Parmi les modules opérationnels livrés figurent la gestion des projets par phases avec
cartographie \textbf{Leaflet}, la production et le suivi des devis, factures et bons de
livraison, le portail client pour la consultation et la validation des documents, la
messagerie interne, la planification des tâches, le suivi des contrats et avenants, ainsi
qu'un module de configuration des rôles et des droits d'accès. Un journal d'activité assure
par ailleurs la traçabilité des opérations sensibles au sein de l'application.

\vspace{0.5cm}
\noindent\textbf{Mots-clés :} BTP, Gestion de projets, Application web, ERP, Laravel, React,
Construction, Devis, Facturation, Planification.
